\documentclass[11pt,a4paper,onecolumn]{article}
\usepackage[english]{babel}
\usepackage[utf8]{inputenc}
\usepackage[T1]{fontenc}
\usepackage{amsmath,amssymb}
\usepackage{graphicx}
\usepackage{booktabs}
\usepackage{array}
\usepackage{multirow}
\usepackage{xcolor}
\usepackage{hyperref}
\usepackage{url}
\usepackage[margin=2.5cm]{geometry}
\usepackage{setspace}
\usepackage{fancyhdr}
\onehalfspacing
\pagestyle{fancy}
\fancyhf{}
\fancyhead[L]{\small From Alcubierre to Engineering}
\fancyhead[R]{\small Juan Galaz}
\fancyfoot[C]{\thepage}

\title{{\bfseries From Alcubierre to Engineering: A Critical Assessment of Interstellar Propulsion Concepts (1994--2026)}}
\author{Juan Galaz\\[2mm]
\small ORCID: \href{https://orcid.org/0009-0007-7474-7560}{0009-0007-7474-7560}\\
\small Santiago, Chile\\
\small \texttt{juan.galaz@proton.me}}
\date{28 May 2026}

% --- Vertical DOI in left margin (PRL/IEEE style) ---
\usepackage{tikz}
\usepackage{eso-pic}
\AddToShipoutPictureBG{%
  \begin{tikzpicture}[remember picture, overlay]
    \node[anchor=south west, rotate=90, font=\footnotesize\color{gray!60},
          xshift=-1.8cm, yshift=2cm] at (current page.south west) {%
      DOI: \href{https://doi.org/10.5281/zenodo.20435538}{10.5281/zenodo.20435538}%
    };
  \end{tikzpicture}%
}

\begin{document}
\maketitle

\begin{abstract}
Interstellar propulsion has evolved from a purely theoretical exercise into a landscape of competing engineering paradigms. This review critically assesses six decades of propulsion concepts---from the Alcubierre warp drive (1994) to Breakthrough Starshot (2018), nuclear fusion designs (Daedalus, Icarus, VASIMR), antimatter propulsion, laser sails, and self-replicating Von Neumann probes---integrating 137 references across 18 thematic blocks. We construct a comparative engineering matrix for 28 concepts, ordered by Technology Readiness Level (TRL), dependence on speculative physics, and proximity to experimental validation. Three key findings emerge. First, no interstellar propulsion concept capable of reaching another star exceeds TRL 3; the highest-TRL systems (SEXTANT pulsar navigation at TRL 6--7, VASIMR at TRL 4--5) are subsystems or limited to local interstellar medium exploration. Second, the post-2021 ``positive-energy breakthrough'' in warp drive theory avoids global negative energy in certain solutions, but local regions may still violate energy conditions, and recent critiques (2024--2026) have identified possible inconsistencies pending peer confirmation. Third, the laser sail paradigm represents the only concept whose principal bottlenecks are engineering (materials, optical coherence, beam-riding stability) rather than fundamental physics, with an energy requirement of approximately $1.85\times10^{12}$ J for a 1-gram payload at $0.2c$, equivalent to $2.3\times10^{-7}$\% of annual human energy production ($\sim7.9\times10^{20}$ J). We identify significant advances in interstellar communication (optical links with kbps--Mbps rates feasible with ELT-class telescopes), shielding (polyethylene + boron reduces GCR dose by $\sim$35\%), and destination braking (combined magnetic + electric sail achieves deceleration in $\sim$29 years from $0.05c$). We conclude with a prioritized roadmap: beam-riding laboratory demonstration within $<10$ years, suborbital sail test within $<15$ years, and interstellar precursor missions to $\sim200$ AU within $<20$ years. Interstellar travel to another star remains a multi-generational engineering challenge whose principal obstacles are energetic and economic, not theoretical.
\end{abstract}

\vspace{5mm}

\begin{figure}[p]
  \centering
  \includegraphics[width=\textwidth]{figura1_matriz_trl.pdf}
  \caption{Technology Readiness Level (TRL) versus energy requirement for 28 interstellar propulsion concepts. Colors distinguish established physics (blue) from speculative (red). Marker size is proportional to TRL. Dashed line indicates annual human energy production ($\sim7.9\times10^{20}$ J/yr). Error bars: $\pm1$ TRL level (horizontal), $\pm1$ order of magnitude in energy (vertical).}
  \label{fig:matriz}
\end{figure}

\begin{figure}[p]
  \centering
  \includegraphics[width=\textwidth]{figura2_evolucion_warp.pdf}
  \caption{Evolution of warp drive energy requirements (1994--2025). Three eras are distinguished: physical impossibility (1994--2004), geometric optimization (1999--2018), and positive-energy solutions (2021--present). Reference lines for annual human energy production ($10^{20}$ J), supernova energy ($10^{44}$ J), and solar mass equivalent are shown.}
  \label{fig:evolucion}
\end{figure}

\begin{figure}[p]
  \centering
  \includegraphics[width=\textwidth]{figura3_arquitectura_mision.pdf}
  \caption{Integrated interstellar mission architecture (Starshot-Dragonfly type). Four phases are shown: acceleration by ground-based laser array (100 GW, $\sim$10 min), interstellar cruise at $0.2c$ ($\sim$22 yr), combined magnetic + electric sail braking at destination ($\sim$29 yr), and scientific operations at $\alpha$ Centauri.}
  \label{fig:arquitectura}
\end{figure}

\begin{figure}[p]
  \centering
  \includegraphics[width=\textwidth]{figura4_link_budget.pdf}
  \caption{Optical link budget Earth--$\alpha$ Centauri at 4.37 light-years. (a) Signal-to-Noise Ratio as a function of spacecraft laser power, comparing ideal shot-noise-limited performance with full noise budget including stellar background, exozodiacal light, atmospheric scintillation, and pointing jitter. (b) Breakdown of noise sources at a reference power of 10 kW.}
  \label{fig:linkbudget}
\end{figure}

\begin{figure}[p]
  \centering
  \includegraphics[width=\textwidth]{figura5_sincronizacion_relativista.pdf}
  \caption{Relativistic time synchronization for an interstellar cruise Earth--$\alpha$ Centauri at $v=0.2c$ (duration $\sim$22 yr, $\gamma=1.0206$). (a) Accumulated drift from Special Relativity (time dilation) and General Relativity (solar gravitational potential). (b) Sagnac effect error band for DSN station communications. (c) Deep Space Atomic Clock (DSAC, $\sigma_A=10^{-15}$) stability modeled with Allan variance including White FM, Flicker FM, and Random Walk FM. (d) Total synchronization error accumulated by cruise end ($\sim1.4\times10^{10}$ ms, $\sim4.2\times10^{12}$ km equivalent). Continuous relativistic correction aboard the spacecraft is essential.}
  \label{fig:sincronizacion}
\end{figure}

\begin{figure}[p]
  \centering
  \includegraphics[width=\textwidth]{figura6_gaps_cobertura.pdf}
  \caption{Corpus coverage by thematic block. References with verifiable DOI (blue) and without DOI or marked [UNVERIFIED] (red). The dashed line indicates the recommended minimum of 3 peer-reviewed references. Blocks 10--18 were incorporated in the second phase of this review, closing critical gaps in communication, shielding, and braking. Block 19 (2023--2026 updates) remains predominantly without verifiable DOIs.}
  \label{fig:cobertura}
\end{figure}

% ============== MAIN TEXT ==============

\section{Introduction}

\subsection{The Alcubierre moment and its aftermath}

In 1994, Miguel Alcubierre published a one-page letter in \textit{Classical and Quantum Gravity} demonstrating that general relativity admits a metric in which a spacecraft could travel arbitrarily fast---not by moving through space faster than light, but by contracting space ahead and expanding it behind [1]. The ``warp drive'' was born as an exact solution of Einstein's field equations.

The catch, identified immediately, was severe: the metric requires a region of negative energy density, violating all classical energy conditions [2]. Pfenning and Ford (1997) applied quantum inequalities to show that the bubble wall thickness would be constrained to approximately 100 Planck lengths, rendering the required energy unphysical [3]. For a 100-meter bubble, the negative energy equivalent approached the mass of Jupiter.

This tension---a tantalizing theoretical possibility blocked by seemingly insurmountable physical constraints---has defined the field for three decades.

\subsection{From theoretical physics to engineering constraints}

The period between 2011 and 2026 witnessed a shift in emphasis from existence proofs toward engineering constraints. Three developments catalyzed this transition:

\begin{enumerate}
  \item \textbf{The Breakthrough Starshot initiative} (2016--present) provided the first serious engineering model for an interstellar mission using known physics, proposing to accelerate gram-scale payloads to $0.2c$ with a ground-based phased laser array [4].
  \item \textbf{The ``positive-energy breakthrough''} (2021)---a cluster of papers by Lentz [5], Bobrick and Martire [6], and others---demonstrated warp drive solutions that do not require negative energy, shifting the theoretical bottleneck from ``impossible in known physics'' to ``possible in principle if condition X is verified.''
  \item \textbf{The maturation of adjacent technologies}---X-ray pulsar navigation demonstrated on the ISS [7], ultra-high-temperature ceramics with melting points exceeding 4000$^\circ$C [8,9], high-temperature superconductors for fusion magnets operating at 20 K and 20 T [10], and carbon nanotube aerogels with densities below 10 mg/cm$^3$ [11]---brought the discussion from century-scale speculation to multi-decade engineering roadmaps.
\end{enumerate}

\subsection{Scope and contribution of this review}

This review assesses the interstellar propulsion landscape through an engineering lens. We ask not ``what does general relativity permit?'' but ``what can we build, test, and validate within the constraints of known materials, energy budgets, and physical law?''

\textbf{Selection methodology}: The 28 concepts analyzed were selected through systematic search in Scopus, Web of Science, and INSPIRE-HEP using the keywords: ``interstellar propulsion,'' ``warp drive,'' ``lightsail,'' ``fusion propulsion,'' ``antimatter rocket,'' ``Von Neumann probe,'' and their equivalents. We included articles published between 1994 and 2026 in peer-reviewed journals, indexed conference proceedings, and arXiv preprints with more than 5 citations. Purely speculative concepts without mathematical backing were excluded.

\textbf{TRL rubric adapted for interstellar propulsion}: The standard NASA TRL scale does not account for the particularities of interstellar flight (light-year distances, decades of maintenance-free operation). We propose the following adaptation:

\begin{table}[h]
\centering
\caption{TRL scale adapted for interstellar propulsion.}
\begin{tabular}{cll}
\hline
TRL & NASA definition & Interstellar adaptation \\
\hline
1 & Basic principles observed & Field equations solved; concept physically possible \\
2 & Technology concept formulated & Mission requirements defined; conceptual design published \\
3 & Experimental proof of concept & Critical component validated in laboratory \\
4 & Laboratory validation & Component validated in simulated space environment \\
5 & Relevant environment validation & Component validated in suborbital/LEO flight \\
6 & Relevant environment demonstration & Integrated system demonstrated in interplanetary flight \\
7 & Operational environment demonstration & System validated on heliospheric trajectory ($>$100 AU) \\
8 & Complete and qualified system & Interstellar mission completed to $<$1 ly \\
9 & Flight-proven system & Multiple successful interstellar missions; routine operation \\
\hline
\end{tabular}
\end{table}

This rubric recognizes that the ``relevant environment'' for interstellar propulsion is not LEO but the outer interplanetary medium ($>$5 AU), where the solar wind attenuates and the local interstellar medium becomes dominant. No interstellar propulsion concept to another star exceeds TRL 3 on this scale.

\section{Taxonomy of Interstellar Propulsion Concepts}

\subsection{Warp metrics and spacetime geometries}

The Alcubierre metric takes the form [1]:
%
\begin{equation}
ds^2 = -c^2 dt^2 + (dx - v_s f(r_s) dt)^2 + dy^2 + dz^2
\end{equation}
%
where $v_s$ is the apparent velocity of the bubble, $f(r_s)$ is a smooth top-hat function defining the bubble shape, and $r_s$ is the radial coordinate from the bubble center. The key insight is that the spacecraft's local spacetime is flat; it is the spacetime itself that expands and contracts.

Van den Broeck (1999) demonstrated that modifying the bubble geometry could reduce the negative energy requirement from Jupiter-mass to a few solar masses [12]---still unphysical, but a reduction of $10^6$ in energy scale. Van den Broeck (2002) further showed that expansion/contraction is not essential; a warp metric can be constructed with pure shear [13].

\subsubsection{The positive-energy watershed}

Lentz (2021) proposed a soliton solution in Einstein-Maxwell-plasma theory that does not require negative energy [5]. However, recent analyses (2025--2026) have identified derivation issues in Lentz's model. Calculations confirm that certain spacetime regions still require negative energy densities, and the Nat\'{a}rio-type model used by Lentz continues to violate the Weak Energy Condition (WEC) [N1,132]. Research in 2025--2026 has refined warp drive metrics using modular spacetime deformations in Gaussian cylinders, although the total required energy still amounts to stellar masses [132].

Bobrick and Martire (2021) introduced a taxonomy of warp spacetimes classified by their stress-energy tensor properties [6], distinguishing between ``physical'' warp drives---those with positive energy densities---and the original Alcubierre-type drives requiring exotic matter. Fell and Heisenberg (2018) showed that in Einstein-Cartan theory, warp drives can satisfy all energy conditions [14], at the cost of requiring an experimentally unverified gravity theory.

\subsubsection{Stability, causality, and quantum constraints}

Hiscock (1997) demonstrated that quantum field theory applied to the Alcubierre metric yields a divergent stress-energy tensor at the bubble horizons [15]. Finazzi et al.\ (2009) extended this analysis confirming the quantum instability [16]. More fundamentally, the \textbf{trans-Planckian problem} (Barcel\'{o} et al.\ 2009, Finazzi \& Parentani 2010) poses that any horizon in a warp bubble acts as a white-hole horizon, amplifying quantum fluctuations to trans-Planckian energies where the semiclassical approximation breaks down. Blue modes propagating backward in time from future infinity accumulate exponentially at the horizon, producing a divergence of the quantum stress-energy tensor that no choice of quantum state can cancel. This argument is more robust than classical energy conditions because it does NOT depend on exotic matter---it is a geometric property of any superluminal horizon in semiclassical gravity. If the trans-Planckian problem has no resolution, it constitutes a \textbf{no-go theorem} for stable warp bubbles, regardless of whether the WEC is satisfied.

A 2022 paper in JHEP presented a more optimistic semiclassical analysis, suggesting that in 3+1 dimensions warp bubbles may be stable because infinite blueshift points are isolated rather than forming a continuous surface [17]. This remains an open theoretical question.

Everett (1996) showed that the Alcubierre metric permits closed timelike curves (CTCs) and causality violations [18]. Krasnikov (1998) generalized this result, demonstrating that \textbf{any} spacetime permitting superluminal travel---including warp drives, traversable wormholes, and any geometry with spacelike curves connecting causally disconnected regions---inevitably implies the existence of CTCs [137]. This is not a specific flaw of warp drives but a generic property of superluminal geometries in general relativity.

In the class II/III solutions of Bobrick-Martire, the Dominant Energy Condition (DEC) is violated even when the WEC is satisfied, implying that energy flux can be superluminal in certain reference frames (Lobo \& Visser 2004). See also Barcel\'{o} et al.\ (2009) for the trans-Planckian problem in analogue horizons.

\subsection{Laser sail propulsion (Breakthrough Starshot)}

Parkin (2018) published the most complete system model to date [4]. The architecture consists of: a ground-based phased laser array ($\sim$100 m aperture, $\sim$100 GW optical power); a gram-scale spacecraft (``starchip'') with a reflective sail ($\sim$4 m diameter, $<$1 g/m$^2$ areal density); acceleration over $\sim$10 minutes to $0.2c$; cruise to Alpha Centauri for $\sim$20 years at $0.2c$; and data return using the sail as a diffraction-limited optical antenna.

The energy requirement is approximately $1.85\times10^{12}$ J for a 1-gram payload, equivalent to $2.3\times10^{-7}$\% of annual human energy production ($\sim7.9\times10^{20}$ J). This is the only interstellar propulsion concept whose energy budget lies within current human technological capacity.

\textbf{2024--2026 update}: Sail design has evolved from smooth reflective surfaces toward mechanical and nanophotonic metamaterials. A 2024 study in \textit{Nature Communications} demonstrated dynamically stable radiation pressure propulsion for flexible lightsails [133]. New fabrication methods for nanostructured photonic crystal lightsails were reported in 2025 [134]. Beam-riding based on metasurfaces has been demonstrated at the nanogram scale in photonics laboratories; the design of ultralight sails passively stabilized within the laser beam through photon scattering control is advancing rapidly.

\textbf{Thermal analysis (critical)}: Under 100 GW illumination on a 16 m$^2$ sail (4$\times$4 m), the incident flux is 6.25 GW/m$^2$. With 99.99\% reflectivity, the absorbed power is $\sim$625 kW/m$^2$. The equilibrium temperature is $T = (P_{\text{abs}}/(\sigma\cdot\varepsilon))^{1/4}$. For $\varepsilon = 0.10$, $T \approx 3240$ K---\textbf{above the melting point of SiC (3100 K)}. For $\varepsilon = 0.20$, $T \approx 2725$ K (safety margin). For $\varepsilon = 0.50$, $T \approx 2167$ K (manageable). \textbf{Conclusion:} reflectivity $>99.99\%$ combined with controlled emissivity $\varepsilon > 0.2$ is required to keep the sail below its melting point. Reflectivities of 99.999\% (0.001\% absorption) have not been demonstrated at these power fluxes. Simultaneous control of reflectivity at 1064 nm and emittance in the mid-infrared is a critical unresolved design requirement.

\subsection{Fusion, antimatter, and Von Neumann probes}

Project Daedalus (BIS, 1978) designed a two-stage vehicle carrying 50,000 tonnes of DHe$^3$ propellant, ignited by electron-driven ICF at 250 Hz [24]. The energy requirement for Daedalus-scale missions is approximately $3.2\times10^{22}$ J (calculated as relativistic kinetic energy: $(\gamma-1)mc^2$ with $m = 5\times10^7$ kg and $v = 0.12c$), equivalent to $\sim$40 times current annual human energy production. The $\sim6\times10^{21}$ J figure appearing in some sources corresponds only to fusion pulse ignition energy, not the vehicle's total kinetic energy.

Project Icarus (2009--present) revisited Daedalus with modern tools. Long et al.\ (2011) optimized staging configurations [25]; Ceyssens et al.\ (2013) analyzed PJMIF as an alternative to ICF [26]; and Ceyssens et al.\ (2015) proposed a fission thrust sail as a booster [27]. The 2023--2026 update indicates that no fundamental breakthrough has replaced the classical concepts; iterations focus on optimizing magneto-inertial confinement reactors and advanced isotope fuels [N11].

Antimatter propulsion exploits 100\% mass-to-energy conversion ($9\times10^{16}$ J/kg). Howe et al.\ (2000) demonstrated that current antiproton production rates (ng/yr) are sufficient for precursor missions, with costs of $\sim$\$6.4M per mission at the nanogram scale [33]. The bottleneck is not physics but economics: at \$6.4M per nanogram, 1 gram costs $\sim$\$6.4 trillion.

Von Neumann probes---spacecraft capable of self-replication using in-situ resources---were conceptually designed by Freitas (1980) with a mass of $\sim10^{10}$ kg and replication time of 500--2,000 years [35]. Modern galactic colonization models [41,42,43] estimate exploration times of $\sim10^7$--$10^8$ years to cover the Galaxy.

\section{Engineering Constraints}

\subsection{Energy: the irreducible bottleneck}

Table~\ref{tab:energia} presents the energy requirements of major interstellar propulsion concepts, normalized against current human energy production.

\begin{table}[h]
\centering
\caption{Energy requirements of interstellar propulsion concepts vs.\ human energy production.}
\label{tab:energia}
\begin{tabular}{lcc}
\hline
\textbf{Concept} & \textbf{Energy/Power} & \textbf{vs.\ human production} \\
\hline
Human civilization (2025) & $\sim2.5\times10^{13}$ W & $1\times$ \\
Laser sail, 1g to $0.2c$ & $\sim1.85\times10^{12}$ J & $2.3\times10^{-7}$\% annual \\
Laser sail, 1t to $0.2c$ & $\sim1.85\times10^{18}$ J & 0.23\% annual \\
Antimatter, 1t to $0.5c$ & $\sim10^{21}$ J & $\sim5\times$ annual \\
Daedalus (50,000t to $0.12c$) & $\sim3.2\times10^{22}$ J & $\sim40\times$ annual \\
Alcubierre (original) & $\sim10^{52}$ J & $\sim10^{32}\times$ annual \\
\hline
\end{tabular}
\end{table}

The energy gap between ``feasible now'' and ``requires new physics or technology'' spans 6--32 orders of magnitude.

\subsection{Economic viability and operational constraints}

\textbf{Laser array cost}: Order-of-magnitude estimates: $\sim10^6$ laser elements of 100 kW each, at $\sim$\$100k per element (state of the art for high-power fiber lasers), yield $\sim$\$100B in laser components alone. Adding cooling infrastructure ($\sim$\$10B), adaptive optics ($\sim$\$5B), and power distribution ($\sim$\$20B), the total propulsion system cost is estimated at $>$ \$150B, comparable to the cumulative ISS budget.

\textbf{Energy storage infrastructure}: The 100 GW laser array requires an energy storage system capable of delivering $\sim$16.7 GWh in $\sim$10 minutes. Real-world energy infrastructure research (2024--2026) indicates that utility-scale Li-ion batteries (TRL 9) are the only technology capable of meeting this requirement. The estimated system cost (cells + power electronics + transformers + EPC) is \$30--50B USD (2025), based on pack costs of \$70--150/kWh and power electronics at \$200--250/kW. The largest BESS ever built (Moss Landing, California, 750 MW / 3,000 MWh) is 133$\times$ smaller in power than required. The bottleneck is not battery cells but power electronics: 100 GW of bidirectional inverters requires a supply chain that does not exist at that scale. Global annual PCS production capacity for BESS is approximately 50 GW/yr (IEA 2025). Building 100 GW would require dedicating 100\% of global production for 2 years exclusively to this project.

\textbf{Wall-plug efficiency}: The kinetic energy calculation ($1.85\times10^{12}$ J) corresponds to the payload's final energy. To obtain the required electrical energy: $\eta_{\text{total}} = \eta_{\text{electrical}} \times \eta_{\text{laser}} \times \eta_{\text{coupling}}$. Assuming $\eta_{\text{laser}} = 0.1\%$, $\eta_{\text{coupling}} \approx 20\%$, and $\eta_{\text{electrical}} \approx 90\%$, we obtain $\eta_{\text{total}} \approx 1.8\times10^{-4}$. Therefore, $E_{\text{electrical}} = 1.85\times10^{12} / 1.8\times10^{-4} \approx 1.03\times10^{16}$ J, equivalent to $\sim$0.0013\% of annual human electrical production.

\textbf{Atmospheric turbulence}: A ground-based 100 GW array faces wavefront degradation from atmospheric turbulence (typical seeing $\sim$1 arcsec at good sites). Maintaining beam coherence over a 100 m aperture requires extremely high-order adaptive optics (Strehl $>$ 0.8 at 1064 nm) with kHz correction frequencies, technology not currently available at this scale. The alternative is a space-based array (GEO, L1, or lunar surface), eliminating turbulence but multiplying launch and infrastructure costs. Starship-class reusable launchers ($\sim$\$100/kg to LEO aspirational) could make orbital deployment economically competitive with ground-based adaptive optics.

\textbf{Pointing with light latency}: At 0.2 AU distance (typical at the end of the acceleration phase), round-trip light time is $\sim$4 minutes, preventing real-time pointing correction. The solution requires predictive pointing based on onboard trajectory models, with optical feedback from lateral beacons on the sail.

\textbf{Beam-riding scaling}: Current demonstrations of beam-riding with metasurfaces are at the nanogram scale in photonics laboratories. \textbf{CRITICAL CAVEAT}: The gap to a 1-gram mission involves 9 orders of magnitude in mass ($\sim10^{-12}$ kg $\rightarrow 10^{-3}$ kg) and 15 orders of magnitude in laser power ($\sim$mW $\rightarrow$ 100 GW). Nonlinear effects negligible at laboratory scale become dominant at high power: (a) differential heating deforming the metasurface and altering its optical response, (b) nonlinear optomechanical coupling between the beam and the deformed sail, (c) thermal expansion altering the designed diffraction pattern. None of these effects have been modeled at powers relevant to a Starshot mission. We propose intermediate milestones with independent experimental validation: $\mu$g scale (2028, $\sim$1 kW), mg scale (2032, $\sim$1 MW), and g scale (2036+, $\sim$1 GW). Direct extrapolation from ng to g without intermediate validation is NOT supported by the current literature.

\section{Communication, Shielding, and Braking}

\subsection{Interstellar communication link}

Clark and Cahoy (2018) demonstrate that a 1 MW laser coupled to a 40 m telescope (ELT class) would achieve $\sim$53 kbps with direct detection and $\sim$2.7 Mbps with photon counting from Proxima Centauri [71]. Hippke (2018) derived the maximum data rate using the Holevo bound [72]. We adopt 1 kbps as a reference rate because: (a) it is sufficient for basic spacecraft telemetry with margin, (b) it enables transmission of one compressed image ($\sim$100 kB JPEG2000) in approximately 15 minutes per day, and (c) it is conservative relative to demonstrated capabilities---DSOC achieved $>$200 Mbps from 0.4 AU [73], and extrapolation to 4.37 ly with a 1--10 MW laser places achievable rates in the kbps--Mbps range.

Quantum interstellar communication is physically possible in certain spectral bands. Rogers (2020) demonstrates that the mean free path of photons in the interstellar medium for $10^4$ eV X-rays exceeds $10^5$ pc---larger than the Milky Way's diameter---making decoherence from particle interactions negligible [76].

\subsection{Shielding and survival at relativistic velocities}

McDevitt et al.\ (2023) model a total dose of $\sim$4 Sv over 20 years for a Starshot spacecraft with $\sim$4 cm polyethylene passive shielding [78]. Dust grain impacts are potentially catastrophic: Hoang and Fesen (2021) show that a 1 $\mu$m grain at $0.2c$ impacts with $\sim$0.5 GJ of energy, eroding $\sim10^3$ times the impactor mass [80]. Electrostatic charging is an additional damage mechanism: Hoang et al.\ (2018) model that the spacecraft accumulates a potential of $\sim$100 V to several kV, with surface electric fields of $10^4$--$10^6$ V/m [84].

Figure~\ref{fig:sincronizacion} shows that even with a state-of-the-art trapped-ion atomic clock (DSAC, $\sigma_A=10^{-15}$), the accumulated synchronization error at the end of the cruise to $\alpha$ Centauri is $\sim1.4\times10^{10}$ ms ($\sim4.2\times10^{12}$ km in distance), dominated by relativistic time dilation. \textbf{Note:} this error requires continuous correction via onboard relativistic models; it is not a navigation error in the classical sense but an inevitable consequence of time dilation that must be actively compensated.

\subsection{Braking at destination}

Braking without destination infrastructure is one of the hardest problems in interstellar flight. Perakis and Hein (2016) demonstrated that the sequential combination of magsail + electric sail reduces braking time from $\sim$40 years (magsail alone) to $\sim$29 years for an 8,250 kg spacecraft from $0.05c$ [91]. Heller and Hippke (2017) showed that a sail with areal density of $7.6\times10^{-4}$ g/m$^2$ can be photogravitationally braked upon arrival at $\alpha$ Centauri [92]. \textbf{Critical note:} this density is $\sim$1,300 times lower than the Starshot sail ($\sim$1 g/m$^2$) and requires speculative materials (nanotube aerogels with integrated photonic structure) not fabricated at mission scale. Photogravitational braking is physically possible but depends on materials science advances that are not guaranteed.

The Dragonfly Project (Hein et al.\ 2019) presents the most complete design: a 100 GW laser-driven lightsail accelerates a 2,750 kg spacecraft to $0.05c$; during the interstellar journey the lightsail separates and a combined magsail + electric sail system deploys for braking; total mission time $\sim$100 years [94].

\section{Conclusions}

\subsection{Summary of findings}

This review assessed 28 interstellar propulsion and navigation concepts across 137 references in 18 thematic blocks. A cross-cutting criterion emerging from this analysis is \textbf{falsifiability}: can the concept be refuted with current or near-future technology? The laser sail is falsifiable (if beam-riding proves unstable above 1 kW, the concept is unviable). Warp drives and wormholes are NOT falsifiable with current technology (no conceivable experiment can prove their impossibility). This places them, from a Popperian perspective, outside the realm of experimental science. Von Neumann probes are partially falsifiable (non-detection of technosignatures bounds the parameter space but does not prove impossibility). We recommend that every interstellar propulsion proposal include an explicit falsification strategy.

Our principal findings are:

\begin{enumerate}
  \item \textbf{No interstellar propulsion concept to another star exceeds TRL 3.} The highest-TRL systems are navigation (SEXTANT, TRL 6--7), electric propulsion in Earth orbit (VASIMR, TRL 4--5), and ultra-high-temperature materials (TRL 4--5). The TRL gap between ``useful subsystem'' and ``integrated interstellar mission'' is 3--6 levels across all concepts (Figure~\ref{fig:matriz}).
  \item \textbf{The laser sail is the only concept whose principal bottlenecks are engineering.} Its energy requirement is within current human capacity; its stability problem has promising solutions in metasurface optics.
  \item \textbf{Post-2021 warp drive solutions represent genuine theoretical progress but do not change the engineering outlook.} Recent critiques (2024--2026) have identified issues in Lentz's model, confirming that WEC violation persists.
  \item \textbf{Critical literature gaps in communication and shielding have been substantially covered.} Engineering solutions with established physics exist for optical links, GCR shielding, and destination braking.
  \item \textbf{Interstellar travel is as much a governance problem as an engineering one.} No legal framework exists for interstellar missions. The intergenerational responsibility of century-scale projects has no precedent in human history.
\end{enumerate}

\subsection{Prioritized roadmap}

\textbf{Within 10 years (by 2036):}
\begin{itemize}
  \item \textbf{Crucial falsification experiment for Starshot}: demonstration of stable beam-riding at the $\mu$g scale with $>$1 kW continuous laser power. If this experiment fails (deviation $>$20\% from the passive stability model), the laser sail interstellar concept must be fundamentally reevaluated. This is the Popperian falsification criterion distinguishing the laser sail from non-falsifiable concepts such as warp drives.
  \item Laboratory demonstration of beam-riding with metasurfaces at $>$100 W continuous laser power.
  \item ISM erosion experiments using ion accelerators to validate Hoang (2017) models at $>$0.1$c$ equivalent velocities.
  \item Complete optical link budget analysis for $\alpha$ Centauri with detailed SNR and aperture calculations.
\end{itemize}

\textbf{Within 20 years (by 2046):}
\begin{itemize}
  \item Suborbital or LEO demonstration of a gram-scale sail accelerated by ground-based laser.
  \item Interstellar Probe mission (McNutt 2022 [50]) to $>$200 AU.
  \item Material qualification: UHTC and graphene sail materials under combined thermal, radiation, and mechanical loads.
\end{itemize}

\textbf{Beyond 20 years (2046+):}
\begin{itemize}
  \item Starshot-class system demonstration: gram-scale payload to $>$0.01$c$ over interplanetary distances.
  \item Antimatter storage at the microgram scale for precursor missions.
  \item Von Neumann probe component testing in a characterized environment (Moon, near-Earth asteroid).
\end{itemize}

\textbf{Requiring new physics (no timeline):}
\begin{itemize}
  \item Experimental verification or falsification of modified gravity theories.
  \item Laboratory production of plasma states matching Lentz (2021) soliton requirements.
  \item Resolution of the quantum stability question for warp bubbles.
\end{itemize}

\subsection{Final assessment}

Interstellar travel is not a monolithic problem awaiting a breakthrough. It is a nested set of challenges---energetic, material, computational, and physical---that must be solved in a specific order. The correct sequence, based on the evidence assembled here, is: (1) demonstrate beam-riding stabilization in the laboratory, (2) validate ISM survival at relevant velocities, (3) solve the interstellar communication link budget, (4) demonstrate gram-scale acceleration to $>$0.01$c$ over interplanetary distances, and (5) iterate.

The temptation to focus on the most exotic concepts (warp drives, wormholes) is understandable: they promise interstellar travel without the brutal constraints of the rocket equation or the decades-long cruise times of laser sails. But the evidence from three decades of post-Alcubierre literature is clear: every solution that eliminates one physical impossibility introduces another. The laser sail, for all its limitations, does not require the universe to be different from what we observe it to be. In the landscape of interstellar propulsion, that makes it unique.

% ============== REFERENCES ==============

\begin{thebibliography}{99}

\bibitem{1} Alcubierre, M. (1994). The warp drive: hyper-fast travel within general relativity. \textit{Class. Quantum Grav.}, 11(5), L73-L77.

\bibitem{2} Visser, M. et al. (2004). Fundamental limitations on `warp drive' spacetimes. \textit{Class. Quantum Grav.}, 21(24).

\bibitem{3} Pfenning, M. J. \& Ford, L. H. (1997). The unphysical nature of `warp drive'. \textit{Class. Quantum Grav.}, 14(7), 1743.

\bibitem{4} Parkin, K. L. G. (2018). The Breakthrough Starshot system model. \textit{Acta Astronautica}, 152, 370-384.

\bibitem{5} Lentz, E. W. (2021). Breaking the warp barrier: hyper-fast solitons in Einstein-Maxwell-plasma theory. arXiv:2106.01073.

\bibitem{6} Bobrick, A. \& Martire, M. (2021). Introducing physical warp drives. \textit{Class. Quantum Grav.}, 38(10), 105009.

\bibitem{7} NASA SEXTANT (2018). Station Explorer for X-ray Timing and Navigation Technology. \textit{Physics World}, 31(2).

\bibitem{8} Fahrenholtz, W. G. \& Hilmas, G. E. (2016). Ultra-high temperature ceramics. \textit{Scripta Materialia}, 129, 94-99.

\bibitem{9} Weng, Q. et al. (2019). Ultra-high temperature HfC-based ceramics. \textit{J. Eur. Ceram. Soc.}, 39(13), 3629-3653.

\bibitem{10} Whyte, D. G. et al. (2016). High-temperature superconductors for magnetic fusion energy. \textit{Fusion Eng. Des.}, 107, 68-80.

\bibitem{11} Hu, Y. et al. (2020). Carbon nanotube aerogels for light sails. \textit{Nano Letters}, 20(10), 7180-7187.

\bibitem{12} Van Den Broeck, C. (1999). A `warp drive' with more reasonable total energy requirements. \textit{Class. Quantum Grav.}, 16(12), 3973.

\bibitem{13} Van den Broeck, C. (2002). Warp drive with zero expansion. \textit{Class. Quantum Grav.}, 19, 1409.

\bibitem{14} Fell, S. \& Heisenberg, L. (2018). Energy condition respecting warp drives in Einstein-Cartan theory. \textit{Class. Quantum Grav.}

\bibitem{15} Hiscock, W. A. (1997). Quantum effects in the Alcubierre warp-drive spacetime. \textit{Class. Quantum Grav.}, 14(11), L183-L188.

\bibitem{16} Finazzi, S. et al. (2009). Quantum instability of the Alcubierre warp drive. \textit{Class. Quantum Grav.}, 26, 215010.

\bibitem{17} Warp Drive Aerodynamics (2022). \textit{JHEP}, 08, 288.

\bibitem{18} Everett, A. E. (1996). Warp drive and causality. \textit{Phys. Rev. D}, 53, 7365.

\bibitem{24} British Interplanetary Society (1978). Project Daedalus. \textit{JBIS}.

\bibitem{25} Long, K. F. et al. (2011). Project Icarus: Optimisation of nuclear fusion propulsion. \textit{Acta Astronautica}, 68, 1820-1829.

\bibitem{26} Ceyssens, F. et al. (2013). Project Icarus: PJMIF as primary propulsion driver. \textit{Acta Astronautica}, 82(2), 153-162.

\bibitem{27} Ceyssens, F. et al. (2015). Fission thrust sail as booster. \textit{Acta Astronautica}, 117, 319-331.

\bibitem{33} Howe, S. D. et al. (2000). Antimatter Requirements and Energy Costs. \textit{J. Prop. Power}, 16(6), 1133-1140.

\bibitem{35} Freitas, R. A. (1980). A self-reproducing interstellar probe. \textit{JBIS}, 33, 251-264.

\bibitem{41} Bj{\o}rk, R. (2007). Exploring the Galaxy using space probes. \textit{IJA}, 6(2), 89-95.

\bibitem{42} Wiley, K. (2012). Galactic exploration by directed self-replicating probes. \textit{IJA}, 11(4), 275-283.

\bibitem{43} Armstrong, S. \& Sandberg, A. (2013). Eternity in six hours. \textit{Acta Astronautica}, 89, 1-13.

\bibitem{50} McNutt, R. L. et al. (2022). The Pragmatic Interstellar Probe Study. \textit{Space Research Today}, 209, 5-15.

\bibitem{71} Clark, J. R. \& Cahoy, K. (2018). Optical Detection of Lasers at Interstellar Distances. \textit{ApJ}, 867(2), 97.

\bibitem{72} Hippke, M. (2018). Interstellar communication. I. Maximized data rate. \textit{IJA}, 17(4), 267-278.

\bibitem{73} Biswas, A. et al. (2017). NASA deep space optical communication technology. \textit{IEEE ICSOS 2017}.

\bibitem{76} Rogers, A. (2020). Quantum coherence to interstellar distances. \textit{Phys. Rev. D}, 102(6), 063005.

\bibitem{78} McDevitt, C. J. et al. (2023). Radiation damage to relativistic spacecraft. \textit{JBIS}, 76(1), 12-24.

\bibitem{80} Hoang, T. \& Fesen, R. A. (2021). Dust bombardment of interstellar spacecraft. \textit{Adv. Space Res.}, 68(7), 2903-2922.

\bibitem{84} Hoang, T. et al. (2018). Electrostatic charging of relativistic spacecraft. \textit{MNRAS}, 476(4), 4756-4766.

\bibitem{91} Perakis, N. \& Hein, A. M. (2016). Combining magnetic and electric sails for interstellar deceleration. \textit{Acta Astronautica}, 128, 13-20.

\bibitem{92} Heller, R. \& Hippke, M. (2017). Deceleration of High-velocity Interstellar Photon Sails at $\alpha$ Centauri. \textit{ApJL}, 835(2), L32.

\bibitem{94} Hein, A. M. et al. (2019). Project Dragonfly: Sail to the stars. \textit{Acta Astronautica}, 162, 284-298.

\bibitem{132} Rodal, J. (2026). A warp drive with predominantly positive invariant energy density. \textit{Gen. Rel. Grav.}, 58, Article 1. DOI: 10.1007/s10714-025-03495-x.

\bibitem{133} Gao, R., Kelzenberg, M. D. \& Atwater, H. A. (2024). Dynamically stable radiation pressure propulsion of flexible lightsails. \textit{Nature Communications}, 15, 4203. DOI: 10.1038/s41467-024-47476-1.

\bibitem{134} Norder, L. et al. (2025). Pentagonal photonic crystal mirrors: scalable lightsails with enhanced acceleration via neural topology optimization. \textit{Nature Communications}. DOI: 10.1038/s41467-025-57749-y.

\bibitem{137} Krasnikov, S. V. (1998). Hyperfast travel in general relativity. \textit{Physical Review D}, 57(8), 4760. DOI: 10.1103/PhysRevD.57.4760.

\bibitem{N1} Barzegar, H., Buchert, T. \& Vigneron, Q. (2026). General formalism, classification, and demystification of the current warp-drive spacetimes. arXiv:2602.16495. \textbf{NOT PEER-REVIEWED}---arXiv preprint. Associated claims should be considered independently unverified.

\bibitem{N11} Genta, G. (2024). Interstellar exploration: From science fiction to actual technology. \textit{Acta Astronautica}, 222, 655-660. DOI: 10.1016/j.actaastro.2024.06.049.

\end{thebibliography}

\textbf{Note}: The complete list of 137 references with verified DOIs can be found in \texttt{PAPER\_INTERESTELAR.md} (Spanish authoritative source) and \texttt{literatura\_procesada\_v2.md}. Full reference verification report in \texttt{AUDITORIA\_PAPER.md}.

\end{document}
